\begin{center}
    {\LARGE \textbf{ĐỀ LUYỆN 04}}\\[6pt]
    {\large (Dành cho đội tuyển quốc gia)}
\end{center}

\begin{tabularx}{\textwidth}{|c|X|c|c|c|}
\hline
\textbf{Bài} & \textbf{Tên bài} & \textbf{File đầu vào} & \textbf{File đầu ra} & \textbf{Điểm} \\
\hline
1 & Chia và nối mảng & \texttt{SADD.INP} & \texttt{SADD.OUT} & 7 \\
\hline
2 & Rắn săn mồi & \texttt{SNAGR.INP} & \texttt{SNAGR.OUT} & 7 \\
\hline
3 & Tạo xâu & \texttt{STRCONS.INP} & \texttt{STRCONS.OUT} & 6 \\
\hline
\end{tabularx}



\section*{Bài 3: Dãy con}
% https://www.codechef.com/problems/CT27

Cho một xâu kí tự $S$ độ dài $N$, và một tập xâu $D$ gồm $M$ xâu kí tự phân biệt. Bạn cần đếm số cách nối các xâu trong $D$ (có thể dùng nhiều lần) để tạo thành xâu $S$.

Ví dụ: $S = \texttt{ABAB}$ và $D = \{\texttt{A}, \texttt{AB}, \texttt{B}\}$. Có 4 cách nối các xâu trong $D$ để tạo thành $S$:

\begin{itemize}
  \item $\texttt{A} + \texttt{B} + \texttt{A} + \texttt{B}$
  \item $\texttt{A} + \texttt{B} + \texttt{AB}$
  \item $\texttt{AB} + \texttt{A} + \texttt{B}$
  \item $\texttt{AB} + \texttt{AB}$
\end{itemize}

\subsubsection*{Dữ liệu vào}

\begin{itemize}
    \item Dòng đầu chứa xâu kí tự $S$ độ dài $N$.
    \item Dòng thứ hai chứa số nguyên $M$ — số lượng xâu trong tập $D$.
    \item M dòng tiếp theo, mỗi dòng chứa một xâu kí tự $D_i$.
\end{itemize}

\subsubsection*{Dữ liệu ra}
In ra số cách nối các xâu trong $D$ để tạo thành xâu $S$. Vì số cách có thể rất lớn, bạn chỉ cần in ra kết quả modulo $10^9 + 7$.

\subsubsection*{Ràng buộc}
\begin{itemize}
  \item $1 \le |S|, M \le 2\times 10^5$
  \item Tổng độ dài các xâu trong $D$ không vượt quá $2\times 10^5$.
\end{itemize}

\subsubsection*{Ví dụ mẫu}

\begin{tcolorbox}[title=Input]
\begin{verbatim}
ABAB
3
A
AB
B
\end{verbatim}
\end{tcolorbox}

\begin{tcolorbox}[title=Output]
\begin{verbatim}
4
\end{verbatim}
\end{tcolorbox}

\textbf{Subtask}

\begin{itemize}
  \item (30\% điểm) $N \le 2000$, tổng độ dài các xâu trong $D$ không vượt quá $2000$.
  \item (70\% điểm) Không có ràng buộc thêm.
\end{itemize}